\section {Self-replicating robotic systems}

\subsection{Self-replication introduction}
Self-replication is a process of construction of an identical or similar copy of itself by some dynamical system.

A self-replicating machine is a type of autonomous robot that is capable of reproducing 
itself autonomously using materials found in the environment, thus exhibiting self-replication 
in a way analogous to that found in nature.

Self-assembly is the autonomous organization of components into patterns or structures
without human intervention. 

In robotics, this can either refer to the ability of a kinematic 
machine to manipulate a series of parts into an assembled copy of itself (kinematic self-replication) or the capability of a
group of mobile robots to autonomously connect to and disconnect from each other (modular
robotics).

This differs from self-replication in that the resulting system is not necessarily
capable of making, catalyzing, or in some way inducing more copies of itself

Difference between replicating system (requires an external control elements for replication e.g. virus)
and Self-replicating system.




\subsection{Classification of Self-replication systems}
 
TODO: add classification from \cite{Tan2020Framework}

\subsubsection{Type-related replication} 
As it shown in dissertation of \citeauthor{Jones2021Thesis} \cite{Jones2021Thesis} and  
works of \citeauthor{Lee2007Robotic} \cite{Lee2007Robotic}, 
self-replicating systems can be classified based on type-related replication into:\\
Centralized replication, SRRS where robots can produce only robots without replication ability,

Decentralized replication, SRRS where robots with replication ability can produce copies of itself.
\begin{figure}[h]
    \centering
    \subfloat[Centralized replication]{ \includegraphics[width=6cm]{ClassificationCentralized.png}}
    \subfloat[Decentralized replication]{ \includegraphics[width=7cm]{ClassificationDecentralized.png}}\\
\end{figure}

Hierarchical replication,  SRRS in which is a combination of centralized and decentralized approaches, 
where robots with ability to replicate can produce both types of robots.


\begin{figure}[h]
    \centering
    \includegraphics[width=6cm]{ClassificationHierarchical.png} 
    \caption{Hierarchical replication}
\end{figure}

\subsubsection{Number of robot types involved:}
Homogeneous - all robot in a system have one type\\ 
Heterogeneous - there are different types of robots\\
 

\subsubsection{Involvement of external system }
\cite{Lee2007Robotic}:\\
Passive self-replicating system - is a system that does not have active functionality for self-replication,
but there exist a replication process associated with that system without
any external constructor. 
Penrose block replicators, can be viewed as passive self-replicating systems \cite{Penrose1958Mechanics}.

Active self-replicating system has sufficient functionality to replicate itself, possibly with aid of some
passive environmental structures.\\

Authors \citeauthor{Suthakorn2002SpaceUtilization} proposed more detailed classification \cite{Suthakorn2002SpaceUtilization}, that describes 
replication behaviour and interaction with external systems. According to this classification all SRRS can be subdivided into:

Directly replicating SRRS - in which a robot capable of producing an exact replica of itself in one generation 

Indirectly replicating SRRS - where robot capable of producing one or more intermediate robots that are in turn
capable of producing replicas of the original.

\begin{figure}[h]
    \centering
    \includegraphics[width=14cm]{ClassificationDirect.png} 
    \caption{Directly and indirectly replicating robots classification}
\end{figure}
Directly Replicating Robots can be further classified into following groups:
\begin{itemize}
    \item Fixture-Based Group
        The self-replicating robots in this group depend on external fixtures in order to complete the self-replication
        process. Passive fixtures are able to assist in this because of the shape constraints that they impose. In some
        other cases, to unify subsystems, push-pull fixtures are helpful as well.
    \item Operating-Subsystem-in-Process Group
        In this group one or several subsystems of the replica can operate before the replica itself is fully
        assembled. These subsystems are able to assist the original self-replicating robot during the assembly of the
        replica.
    \item  Single-Robot-Without-Fixture Group
        In this group only one robot is used to finish the self-replication process. Thus, the robot in this group
        depends only on the available environment. Usually, the complexity of the subsystems or the number of
        subsystems in the replica is very low for this group.
    \item  Multi-Robot-Without-Fixture Group
        In this group more than one robot works together in the self-replication process without the assistance of
        fixtures. A major advantage is the reduction of the time required for self-replication. A disadvantage is that
        there may be interference problems among robots.
        There are several possible ways that a self-replicating robot can be categorized in two or three groups
        mentioned above. The combination of two or three different concepts can be incorporated in a potential
        design, such as a combining operating-subsystems-in-process with fixture-based robots. More categories
        are likely to be developed in the subsequent stages of our research in the area of self-replicating robots.
\end{itemize}

Indirectly Replicating Robots or group of robots works together build a robot-producing factory 
or some type of intermediate robot which is able to produce replicas of the original robot.




\subsection{Self-replication models}



First self-reproducing automation model was proposed by \citeauthor{vonNeumann1962Theory} in \citeyear{vonNeumann1962Theory}  \cite{vonNeumann1962Theory}\\
According to this model, a self-reproducing automaton consists of four components:\\
A - an automatic factory (constructor) that collects raw materials and processes them into outputs using a
specified written instruction, \\
B - a copier that takes an instruction and copies it,\\
C - a controller that controls A and B, actuating them alternatively,\\
D - a list of written instructions.\\
\[ (A + B + C + D) \rightarrow (A + B + C + D), (A + B + C + D) \]
And the replication process in von Neumann’s model, is when the automaton 
$ (A + B + C + D) $ produces another $ (A + B + C + D) $.\\
This model in may be incomplete when applied to physical self-replicating systems since it doesn't describe 
interaction with environment, resources and possible leftovers.

\citeauthor{Lee2007Robotic} proposed broader model, where a replication process can
be written as
\[ (R, M, E) \rightarrow (R' , M' , E' ) \]

where: \\
$M$ is a multiset of available parts to be used for building replicas by an initial system \\
$R$ is a multiset of an initial functional system to be reproduced. $R \neq 0$; to be a replicating system. \\ 
$E$ is a multiset of environmental structures and/or elements involved in the replication process (but not 
replicated). \\
Here R describes the whole replicator $ A+B+C+D $  from Von Neumann's model.

Or as a function: 
\[ (R' , M' , E' ) = \Phi(R, M, E, t) \],
where t - time.

Environmental structures, E, can be divided into three categories \cite{Eno2007ReplicationNoControl}:\\
\begin{itemize}

    \item Completely structured environment - in which every environmental structure is fixed at a precise location 
        and there is no modification or change in any environmental
        structures during the self-replication process. 
        If we define $E'$ is the set of environmental structures after self-replication process:
        $E = E' \neq \phi $ 
    \item Partially structured environment is like a structured environment, but
        with structures whose locations can be permuted and perturbed 
        by a small random error in pose during the self-replication process.
        $ E \neq E \neq \phi $ 
    \item Unstructured environment has no predefined or organized structure and 
        can dynamically change (with some independent process $\phi$) during the self-replication.
        Can be formally depicted as:  
        $ E = \phi $ 
\end{itemize}


\subsection{Degree of replication}
In \citeyear{Lee2007Robotic} \citeauthor{Lee2007Robotic} proposed a measure of degree of replication \cite{Lee2007Robotic}:
measure of how much order is created by the assembly process relative to the existing order in
the unassembled parts. 
In this section, we first define a combined measure 
of system complexity distribution and relative complexity for ARS.

As a simple measure of subsystem complexity, we count the number of active elements in
each subsystem. An active element is a moving mechanical part or a fundamental electronic component.
Some special parts, such as a chassis and fixed mechanical parts with special purpose (e.g., a passive end-
effector), are also viewed as active elements even if they are not moving mechanical parts.


For a robotic system (recall that all robotic systems are active by our definition) 
consisting of n subsystems, we define the degree of replication for the system, $D_s$ , as
\begin{equation}\label{eq:1}
    D_s = \frac{C_{min}}{C_{max}} \cdot \frac{C_{total}}{C_{ave}} \cdot \frac{1}{C_{ave}} 
\end{equation}
where
\[ C_{ave} = \frac{1}{n} \sum_{i=1}^{n}{C_i} \] 
\[ C_{max} = max{C_1,...,C_n} \]
\[ C_{min} = min{C_1,...,C_n} \]
\[ C_{total} = \sum_{i=1}^{n}{C_i} + \sum_{i=1}^{n-1} \sum_{j=i+1}^{n}{C_{ij}} \]

$C_i$ is the number of active elements in $M_i$ \\ 
$C_{ij}$ is the number of interconnections between the i-th and j-th subsystems when they are assembled. 
$C_{ij}= 0$ if the i-th and j-th subsystems are not directly connected by the assembly process. 

In formula \ref{eq:1} the first term measures the complexity distribution
throughout the subsystems, the second term measures the relative 
complexity of the total system to the individual subsystems, 
and the last term penalizes for complex subsystems. 
To design a robotic system capable of replication from low-level
parts, the subsystem complexity should remain low relative to the
total system complexity (i.e., high relative complexity). 

In addition, we view a system consisting of a few complex parts as more trivial
than one composed of many low-complexity parts. A higher value
of $D_s$ indicates a more truly replicating system in terms of its 
complexity distribution and relative complexity. Although the concept
of $D_s$ may not generalize to all robotic systems, it is useful as a
measure to evaluate robotic replicating systems.



\subsection{Material-robot systems}
Material-robot systems \cite{AbdelRahman2022Hierarchical}
